\documentclass[a4paper,11pt]{article}

\usepackage{fontspec}
\setmainfont{Carlito}
\usepackage[margin=0.6in, top=0.5in, bottom=0.5in]{geometry}
\usepackage{titlesec}
\usepackage{enumitem}
\usepackage{hyperref}
\usepackage{xcolor}
\usepackage{microtype}
\usepackage{needspace}

\hyphenpenalty=10000 \exhyphenpenalty=10000 \emergencystretch=3em

\definecolor{headernavy}{RGB}{12, 83, 166}
\definecolor{linkblue}{RGB}{0, 102, 204}
\definecolor{muted}{RGB}{96, 102, 110}

\hypersetup{colorlinks=true, linkcolor=linkblue, urlcolor=linkblue}

\titleformat{\section}
  {\color{headernavy}\normalsize\bfseries}
  {}{0em}{\MakeUppercase}[\color{headernavy}\titlerule]

\titlespacing{\section}{0pt}{8pt}{4pt}

\setlist[itemize]{noitemsep, topsep=1pt, leftmargin=1.3em, partopsep=0pt, parsep=1pt}

\pagestyle{empty}
\setlength{\parindent}{0pt}
\linespread{1.02}

% ---- building blocks: bold = organisation/project, italic = role/stack, grey = dates/places ----
% v3.5.3 (10 Nov 2026): employer (and place) on line 1, the title then the dates on line 2, so that `pdftotext -raw` reads every role as
% employer -> title -> dates -> bullets (ats_structure.role_order_problems). The dates were on line 1 before, read before the title.
\newcommand{\role}[4]{% org, location, dates, title
  \textbf{#1}{\color{muted}\,\textperiodcentered\,\textit{#2}}\\
  \textit{#4}\hfill{\color{muted}#3}\par}
\newcommand{\roleNoLoc}[3]{\textbf{#1}\\ \textit{#3}\hfill{\color{muted}#2}\par}
\newcommand{\project}[3]{% name, stack, right-hand tag
  \textbf{#1}\hfill#3\\ {\color{muted}\small\textit{#2}}\par}
\newcommand{\gh}[1]{\href{#1}{\small GitHub}}
\newcommand{\badge}[1]{{\color{headernavy}\small\textbf{#1}}}
\newcommand{\skillgroup}[2]{{\color{headernavy}\textbf{#1}}\\ #2\par\vspace{3pt}}
\newcommand{\gap}{\vspace{5pt}}
\newcommand{\num}[1]{\textbf{#1}}   % result numbers only


\begin{document}

% ================= HEADER =================
\begin{center}
  {\LARGE\bfseries Thiruvel Andagurunathan Pandian} \\[3pt]
  {\color{headernavy}\small\textbf{AI Engineer | LLM agent with deterministic hallucination guards, \textasciitilde{}65 tests}} \\[3pt]
  {\small
  Bristol, UK\nobreak\ \,\textbar\, \href{mailto:apthiruvel@gmail.com}{apthiruvel@gmail.com}\nobreak\ \,\textbar\, \href{tel:+447741016954}{+44 7741 016954}\nobreak\ \,\textbar\, \href{https://www.linkedin.com/in/thiruvel-a-p/}{https://www.linkedin.com/in/thiruvel-a-p/}\nobreak\ \,\textbar\, \href{https://github.com/Thiruvel-AP}{https://github.com/Thiruvel-AP}\nobreak\ \,\textbar\, \href{https://thiruvel-ap.github.io/PersonalPortfolio/}{https://thiruvel-ap.github.io/PersonalPortfolio/}}\par\vspace{4pt}
  \noindent{} Right to Work in UK \,\textbar\, Available to join Immediately
\end{center}

% ================= PROFILE =================
\needspace{5\baselineskip}\section{Profile}
AI engineer who makes AI answers trustworthy: for my AstraZeneca-affiliated MSc dissertation I built an assistant that checks its answers against official databases and keeps working when they go down, checked by about 65 tests. Deploys on AWS (Lambda) with CI/CD and MLflow experiment tracking. Builds the model a problem needs and says plainly where it falls short: customer-churn prediction, price forecasting, text grouping, and models for 11,459 genes that scored 0.807 on training data but 0.545 on unseen data. At 180 Degrees Consulting, led the design in a five-person team that turned a startup\textquotesingle{}s vague idea into a working AI prototype in three months (rated 5 stars); it is now in live use.

% ================= SKILLS =================
\needspace{5\baselineskip}\section{Skills}
\skillgroup{Agentic AI \& LLM-based Applications}{Tool use and tool calling with failure recovery (LangChain, LangGraph, Google ADK, ReAct agent loops; circuit breakers, fallbacks, deterministic guardrails), multi-provider LLMs (OpenAI, Anthropic Claude on AWS Bedrock, Google Gemini via LiteLLM, open-weight models via Ollama), LLM evaluation and testing (three-tier tests, hallucination-rate analysis), agent memory (persistent storage and retrieval across sessions), RAG pipelines (retrieval, generation, embeddings and the ChromaDB vector database), MCP servers (FastMCP, TypeScript SDK), prompt engineering, AI-assisted development (Claude Code as primary tool, Cursor, OpenAI Codex)}
\skillgroup{Software Engineering \& Integrations}{TypeScript/Node.js, React, Python, JavaScript, FastAPI, Pydantic, RESTful APIs and third-party API integrations, GraphQL, PostgreSQL, SQL, MongoDB and DynamoDB (NoSQL), graph data (networkx), pytest, version control with Git (pull requests, code review), debugging and root-cause analysis, stakeholder communication (live client demos, pull-request reviews with a client CTO)}
\skillgroup{Cloud \& Deployment}{AWS (Bedrock, Lambda, ECR, S3, SQS, EC2, IAM, Cognito), Google OAuth sign-in, Docker, Kubernetes (AWS EKS), GitHub Actions CI/CD, serverless deployment}

% ================= EXPERIENCE =================
\needspace{5\baselineskip}\section{Work Experience}
\needspace{4\baselineskip}
\role{AstraZeneca}{Bristol, UK (Remote)}{May 2026 – Sep 2026}{Data Scientist, MSc Dissertation Project}
\begin{itemize}
  \item Built an AI assistant researchers could trust: it checks gene names against official databases and falls back to a local list of 19,213 genes if they go down (LangChain, AWS).
  \item Integrated the agent into a React front end with real-time gene alias enrichment.
  \item Managed AWS account governance for the project, including root-user isolation, IAM user and group management and least-privilege policies.
  \item Found where the assistant went wrong with about 65 automated checks at three levels, fixed failures layer by layer, and reused the checks after changing AI provider (pytest).
  \item Combined 14 separate research datasets so the assistant only used trusted data, then put it online on AWS and documented how it works (Python, AWS Lambda).
\end{itemize}
\gap
\needspace{4\baselineskip}
\role{180 Degrees Consulting Bristol}{Bristol, UK (Remote)}{Sep 2025 – Dec 2025}{Data Consultant (AI Engineering \& AI System)}
\begin{itemize}
  \item Led the design in a five-person team that turned a startup client\textquotesingle{}s vague idea into a working AI prototype in three months; senior leadership rated it 5 stars (Python, MCP).
  \item Proposed the request flow: a user\textquotesingle{}s text goes through the API and a customised prompt to the MCP server, where one LLM instance writes GraphQL queries against the client\textquotesingle{}s Convex database and a second instance runs the statistical analysis before the result returns to the user.
  \item Helped the client\textquotesingle{}s engineers put the tool into everyday use and reviewed their code changes with the CTO; it is \num{now in live use} (TypeScript, OpenAI).
  \item The client\textquotesingle{}s founder later recommended me in writing, calling the server \textquotesingle{}stable, production-ready\textquotesingle{} and saying it arrived ahead of schedule.
\end{itemize}
\gap
\needspace{4\baselineskip}
\role{Freelance}{Chennai, India}{Jan 2025 – Aug 2025}{AI/ML Engineer}
\begin{itemize}
  \item Built a private client an AI finance assistant that reads his income, spending, loans and investments and suggests next steps, self-hosted to keep costs low (Llama, Ollama).
  \item Designed reproducible scikit-learn pipelines with feature engineering (temporal features, interaction terms, target encoding), model selection, cross-validation and hyperparameter tuning across classification and regression tasks.
  \item Working independently for a private client, built agentic RAG workflows (LangChain, LangGraph, ChromaDB) self-hosted on the client\textquotesingle{}s own server, NLP pipelines (Hugging Face Transformers, Ollama) and scikit-learn pipelines, served ML models as REST APIs (FastAPI, Docker), and contributed to a team LoRA fine-tune of a Llama model.
\end{itemize}
\gap
\needspace{4\baselineskip}
\role{Avasoft}{Chennai, India}{Jan 2023 – Sep 2024}{Software Development Engineer Trainee, then Software Development Engineer}
\begin{itemize}
  \item Implemented offline-first data persistence (Isar) with background synchronisation in the Flutter FinTech app, keeping financial records available in low-connectivity conditions.
  \item Helped build a money-management app for a company\textquotesingle{}s 10,000+ employees, linking live stock prices and bank accounts to a backend of 200+ endpoints; shipped the first release (Flutter, AWS).
  \item \num{Shipped features in a production fintech (financial-services) application} as part of a team: integrated Twelve Data (real-time market feeds) and Plaid (bank linking, transaction imports) via REST APIs, with SSL pinning and AWS Cognito SSO.
\end{itemize}

% ================= PROJECTS =================
\needspace{5\baselineskip}\section{Projects}
\needspace{4\baselineskip}
\setbox0=\hbox{\textbf{GTAI Cell Line Finder: multi-omics graph for cancer cell-line selection (PoC)}\hspace{1em}\url{https://github.com/Thiruvel-AP/GTAI-GraphPOC}}%
\setbox2=\hbox{{\small\textit{Python, networkx, Leiden, scikit-learn, RAPIDS cuML}}\hspace{1em}{\small\url{https://github.com/Thiruvel-AP/GTAI-GraphPOC}}}%
\ifdim\wd0>\linewidth
\ifdim\wd2>\linewidth
\textbf{GTAI Cell Line Finder: multi-omics graph for cancer cell-line selection (PoC)}\\ {\color{muted}\small\textit{Python, networkx, Leiden, scikit-learn, RAPIDS cuML}}\\ \hspace*{\fill}{\small\url{https://github.com/Thiruvel-AP/GTAI-GraphPOC}}\par
\else
\textbf{GTAI Cell Line Finder: multi-omics graph for cancer cell-line selection (PoC)}\\ {\color{muted}\small\textit{Python, networkx, Leiden, scikit-learn, RAPIDS cuML}}\hfill{\small\url{https://github.com/Thiruvel-AP/GTAI-GraphPOC}}\par
\fi
\else
\project{GTAI Cell Line Finder: multi-omics graph for cancer cell-line selection (PoC)}{Python, networkx, Leiden, scikit-learn, RAPIDS cuML}{\url{https://github.com/Thiruvel-AP/GTAI-GraphPOC}}
\fi
\begin{itemize}
  \item Combined mRNA and protein evidence per gene with Stouffer\textquotesingle{}s method.
  \item Linked 14 public cancer-research files into one network of 1,399 cell lines and 16,385 genes; models scored 0.807 on training data but 0.545 on unseen data, reported openly (Python).
\end{itemize}
\gap
\needspace{4\baselineskip}
\setbox0=\hbox{\textbf{llama3-8b-qlora-math-t4: QLoRA fine-tuning of Llama 3.1 8B for maths under a token budget}\hspace{1em}\url{https://github.com/Thiruvel-AP/llama3-8b-qlora-math-t4}}%
\setbox2=\hbox{{\small\textit{Python, Unsloth, PEFT, TRL, Hugging Face, PyTorch, sentence-transformers, scikit-learn}}\hspace{1em}{\small\url{https://github.com/Thiruvel-AP/llama3-8b-qlora-math-t4}}}%
\ifdim\wd0>\linewidth
\ifdim\wd2>\linewidth
\textbf{llama3-8b-qlora-math-t4: QLoRA fine-tuning of Llama 3.1 8B for maths under a token budget}\\ {\color{muted}\small\textit{Python, Unsloth, PEFT, TRL, Hugging Face, PyTorch, sentence-transformers, scikit-learn}}\\ \hspace*{\fill}{\small\url{https://github.com/Thiruvel-AP/llama3-8b-qlora-math-t4}}\par
\else
\textbf{llama3-8b-qlora-math-t4: QLoRA fine-tuning of Llama 3.1 8B for maths under a token budget}\\ {\color{muted}\small\textit{Python, Unsloth, PEFT, TRL, Hugging Face, PyTorch, sentence-transformers, scikit-learn}}\hfill{\small\url{https://github.com/Thiruvel-AP/llama3-8b-qlora-math-t4}}\par
\fi
\else
\project{llama3-8b-qlora-math-t4: QLoRA fine-tuning of Llama 3.1 8B for maths under a token budget}{Python, Unsloth, PEFT, TRL, Hugging Face, PyTorch, sentence-transformers, scikit-learn}{\url{https://github.com/Thiruvel-AP/llama3-8b-qlora-math-t4}}
\fi
\begin{itemize}
  \item Fine-tuned an 8-billion-parameter AI model on a free GPU, training only 0.52\% of it, and raised correct maths answers from 16.2\% to 27.3\% on unseen problems (QLoRA, Llama).
  \item Checked all 859,494 training examples for gaps and copies, found 44,369 extra duplicates and hand-reviewed 181 suspect records before training (Python).
\end{itemize}
\gap
\needspace{4\baselineskip}
\setbox0=\hbox{\textbf{AgenticFriend: real-time voice assistant}\hspace{1em}\url{https://github.com/Thiruvel-AP/adk_multi-agent}}%
\setbox2=\hbox{{\small\textit{FastAPI, Google ADK, React, faster-whisper, PyTorch, Docker/CUDA}}\hspace{1em}{\small\url{https://github.com/Thiruvel-AP/adk_multi-agent}}}%
\ifdim\wd0>\linewidth
\ifdim\wd2>\linewidth
\textbf{AgenticFriend: real-time voice assistant}\\ {\color{muted}\small\textit{FastAPI, Google ADK, React, faster-whisper, PyTorch, Docker/CUDA}}\\ \hspace*{\fill}{\small\url{https://github.com/Thiruvel-AP/adk_multi-agent}}\par
\else
\textbf{AgenticFriend: real-time voice assistant}\\ {\color{muted}\small\textit{FastAPI, Google ADK, React, faster-whisper, PyTorch, Docker/CUDA}}\hfill{\small\url{https://github.com/Thiruvel-AP/adk_multi-agent}}\par
\fi
\else
\project{AgenticFriend: real-time voice assistant}{FastAPI, Google ADK, React, faster-whisper, PyTorch, Docker/CUDA}{\url{https://github.com/Thiruvel-AP/adk_multi-agent}}
\fi
\begin{itemize}
  \item Built a voice assistant you can interrupt mid-sentence, running its own speech models on a GPU to cut waiting time (FastAPI, React).
\end{itemize}

% ================= EDUCATION =================
\needspace{5\baselineskip}\section{Education}
\noindent\textit{MSc Data Science}, \textbf{University of Bristol, UK}\unskip\nobreak\hfill\penalty50\hskip1em\null\nobreak\hfill\mbox{Sep 2025 – Sep 2026}\par
\noindent{\small Dissertation with AstraZeneca affiliation (GeneTraceAI)}\par
\noindent\textit{B.E. Electrical \& Electronics Engineering}, \textbf{St. Joseph\textquotesingle{}s College of Engineering}, First Class with Distinction\unskip\nobreak\hfill\penalty50\hskip1em\null\nobreak\hfill\mbox{Aug 2019 – Apr 2023}\par

% ================= CERTIFICATIONS & ACHIEVEMENTS =================
\needspace{5\baselineskip}\section{Certifications \& Achievements}
\begin{itemize}
  \item Model Context Protocol: Advanced Topics, Anthropic, \textit{Jun 2026}
  \item Introduction to Model Context Protocol, Anthropic, \textit{Jun 2026}
  \item AWS Academy Cloud Foundations, \textit{Nov 2025}
  \item 180 Degrees Consulting Data Science Consultant Certificate, \textit{Dec 2025}
  \item Bristol PLUS Award, \textit{Mar 2026}
\end{itemize}

\end{document}